\section{Substitution growth and shift equivalence}

\begin{theorem}[Cross-dimension shift equivalence]\label{thm:cross-dimension-shift-equivalence}\notready\uses{thm:universality-se-class, thm:universality-cross-dim-se}
\emph{(Open frontier --- stated, not formalized.)} \texttt{Universality.lean} scopes to fixed matrix
dimension; the general shift-equivalence claim can change rank/alphabet size. The
target: shift equivalence of primitive substitution systems defined across different alphabet sizes,
with the resulting universality relation still preserving the cohomological classifying invariant
through the moduli stack. Missing bridge: rectangular matrices / dimension groups / alphabet-change
maps. (Target: \texttt{Obsmetrologia.AperiodicOrder.Tiling.Universality.crossDimensionSE}.)
\end{theorem}

\begin{theorem}[Step-$n$ data co-increments]\label{thm:stepn-co-increment}\lean{Obsmetrologia.AperiodicOrder.Tiling.StepN.step_n_co_increments}\leanok\uses{thm:clinv-dichotomy}
For the canonical step-$n$ rule, matrix dimension, recurrence depth,
alphabet size, and cohomological rank are one parameter viewed four ways.
Passing from $n$ to $n+1$ increments all four together.
\end{theorem}

\begin{theorem}[Step-$n$ characteristic polynomial has degree $n$]\label{thm:stepn-growth-degree}\lean{Obsmetrologia.AperiodicOrder.Tiling.StepN.substitution_growth_degree_n}\leanok\uses{thm:stepn-co-increment}
The minimal polynomial of a primitive step-$n$ substitution's growth rate is monic of degree
exactly $n$: Fibonacci ($n=2$) gives $x^2-x-1$ (the golden ratio $\varphi$), Padovan ($n=3$) gives
$x^3-x-1$ (the plastic number $\rho$); the cohomological rank co-increments with $n$, and at $n=2$
the step-$n$ invariant aligns with FOUR (\texttt{substitution\_growth\_degree\_n}).
\end{theorem}

\begin{theorem}[Step-$n$ Perron growth rate]\label{thm:stepn-perron-growth}\notready\uses{thm:stepn-growth-degree}
\emph{(Open frontier --- stated, not formalized.)} The growth rate itself: the dominant
Perron--Frobenius eigenvalue of the $n\times n$ substitution matrix --- the real degree-$n$
algebraic number that is the largest root of the monic characteristic polynomial
(\cref{thm:stepn-growth-degree}), $\varphi$ for Fibonacci and $\rho$ for Padovan. Stating it needs
the reals and the Perron--Frobenius existence/uniqueness theorem, absent from the core. Red until a
constructive real field and Perron--Frobenius are available.
\end{theorem}

\begin{theorem}[The substitution matrix is the Chu $r$ over FOUR, and satisfies its characteristic polynomial]\label{thm:stepn-chu-cayley-hamilton}\lean{Obsmetrologia.AperiodicOrder.Tiling.StepN.cayleyHamilton_padovan}\leanok\uses{thm:stepn-growth-degree}
\cref{thm:stepn-growth-degree} carried the characteristic polynomials as \emph{free-standing}
coefficient lists (\texttt{charPolyPad}$=[-1,-1,0,1]$, \texttt{charPolyFib}$=[-1,-1,1]$) --- monic,
of the right degree, but with no matrix they were the polynomial \emph{of}. This node supplies the
matrix and the connection, entirely from existing pieces (\texttt{Universality.Mat}/\texttt{matMul}/
\texttt{matPow}, \texttt{LinearLogicChu.ChuSpace}, \texttt{StepN.charPolyPad}). The substitution
matrix is the documented abelianization \texttt{substMatrix} (entry $i\,j$ = count of letter $j$ in
the image of $i$); for Padovan it is $[[0,1,0],[0,0,1],[1,1,0]]$. \textbf{This matrix is the $r$ of
a Chu space valued in FOUR}: \texttt{padovanMatrix\_is\_chu\_incidence} proves, term-mode and
\emph{per cell by \texttt{rfl}}, that \texttt{padovanMatrix} is the \texttt{pos}-incidence of
\texttt{padovanChu.r} ($\mathbf t$ present, $\bot$ absent) --- a definitional identity, not a
checked coincidence (step-$n$ land is Chu-$r$ land; at $n=2$ the space \emph{is} FOUR, via
\texttt{step\_two\_is\_four}). \textbf{Cayley--Hamilton} (\texttt{cayleyHamilton\_padovan},
\texttt{cayleyHamilton\_fibonacci}): \texttt{evalMatPoly charPolyPad padovanMatrix} $=0$, i.e.
$M^3=M+I$ --- the plastic number's defining equation $x^3=x+1$ realized on the matrix; likewise
$M^2=M+I$ for Fibonacci. \texttt{padovan\_not\_quadratic} confirms it is genuinely cubic
($M^2\neq M+I$). Proofs are term-mode structural \texttt{Fin}-matches (no \texttt{decide}); the
eigenvalue itself stays the open \cref{thm:stepn-perron-growth} (reals exist in \texttt{Obsmetrologia/Calculus/Real/}
--- \texttt{MetallicReal}/\texttt{Salem} realize specific eigenvalues; the general step-$n$ PF growth
rate across all families is what remains open, not reals), and the
integer-sequence inequality $\rho<\varphi$ is \cref{thm:collatz-padovan-slow-growth}.
\end{theorem}

\begin{theorem}[Cyclic stage maps satisfy $x^N - 1$: the odometer matrix family]\label{thm:permutation-cayley-hamilton}\lean{Obsmetrologia.AperiodicOrder.Tiling.StepN.cayleyHamilton_vershik}\leanok\uses{thm:stepn-chu-cayley-hamilton, thm:vershik-serial-chart}
Cayley--Hamilton as a \emph{family} theorem, not a fixed-size \texttt{rfl}
check: permutation matrices multiply by composing maps
(\texttt{matMul\_permMatrix}) and power by iterating
(\texttt{matPow\_permMatrix}), so any self-map of $\mathrm{Fin}\,n$ whose
$N$-fold iterate is the identity has permutation matrix annihilated by the
monic degree-$N$ list \texttt{charPolyCycle}$\,N = x^N - 1$
(\texttt{cayleyHamilton\_of\_cycle}). The dyadic odometer instantiates it at
every size: \texttt{cayleyHamilton\_vershik} annihilates the stage-$n$
successor's matrix by $x^{2^n} - 1$, the annihilating polynomial realized
for an infinite matrix family from the stage cycle
(\cref{thm:vershik-serial-chart}), where the Fibonacci/Padovan instances are
per-cell at fixed size.
\end{theorem}
